\documentclass[a4paper]{article}

\begin{document}

\noindent o2343003 \\
ONGONWOU FREDERIC\\

\vspace{30pt}


The paper 
\textbf{"Localized nonlinear dissipative matter waves controlled by 
quantum fluctuations"} address the interesting problem of solitons 
in Bose Einstein condensates in an effort to shed light on both.
This is a tempting reserach direction given the similarities between 
solitons and Bose Einstein condensates, both generally described 
by some form of the nonlinear schr\"odinger equation, the latter hodling
generally hodling an additional term wichi is the influence 
of some external field. \\
In this study, the authors go beyond the well known mean field approximation 
and study the dynamics of the system up to the 4th order in an effort to better 
take into account the effects of the quantum fluctuations in the system. 

\section{Theory}

The basic expression to address BEc is the following Gross-Pitaevskii equation
\begin{equation}
i\hbar\frac{\partial \psi}{\partial t}
  = \left[ -\frac{\hbar^2}{2m}\frac{\partial^2}{\partial x^2} + V_{\rm ext}(x) 
+ g\left|\psi \right|^2
\right]\psi, \label{eqn:grosspit}
\end{equation}

\noindent where $g$ is the interaction strength and $V_{\rm ext}$ is 
the external trapping potential, usually in an harmonic form.
Let us note that the presence of $V_{\rm ext}(x)$ is the main difference 
between the above Gross Pitaevskii equation and the
nonlinear Schr\"odinger equation below used to describe solitons in general :
\begin{equation}
i\hbar\frac{\partial \psi}{\partial t} =
   \left[ -\frac{1}{2}\frac{\partial^2}{\partial x^2}
 + g\left|\psi \right|^2\right]\psi.
\end{equation}

The harmonic trap in the Bose Einstein condensate (\ref{eqn:grosspit})
ensures the localization of the matter wave in finite region of space. \\

In this study however, the authors go beyond the mean field approximation 
expressed by the neglection of the higher order term and go up he 4th order

\begin{equation}
i\hbar\frac{\partial \psi}{\partial t}
  = \left[ -\frac{\hbar^2}{2m}\frac{\partial^2}{\partial x^2} + V_{\rm ext}(x) 
+ g_2\left|\psi \right|^2
+ g_3\left|\psi \right|^3
+ g_4\left|\psi \right|^4
 + i\frac{\gamma}{2}
\right]\psi, \label{eqn:main_eq}
\end{equation}

The additional terms help modelize the nature of their specific system
which takes into account :

\begin{itemize}
\item the quantum fluctuations $\sim\left| \psi\right|^3$
\item the three body loss effect $\sim\left| \psi\right|^4$
\item the feeding parameter $\sim\gamma$
\end{itemize}

The precise details which further charaterize their system 
surch as the correction to the external potential caused by the cheosen 
optical lattices are not explicited 
so we can keep it concise.

\section{Method}

The authors then proceed to solve (\ref{eqn:main_eq}) by mean of 
the variational approach, using a classical gaussian trial function

\begin{equation}
\psi(c,\tau) = A(\tau)
\exp\left\lbrace -\frac{c^2}{2a^2(\tau)}
 + \frac{i}{2}b(\tau)c^2  + i\theta(\tau)\right\rbrace
\end{equation}

in the Lagrangian of the system.

Inserting the previous trial function in bthe Lagrangian of the system, 
they are left with the
differential equations for the parameters $A$, $a$ and $b$
which are then solved and analyzed for stability, by the means of 
the Vahkitov-Kolokolov (VK) condition, that states that one of the condition 
for the stability of the soliton is $\partial N/\partial \mu < 0$ everywhere,
where $N$ is the number of particles and $\mu$ the chemical potential

\section{Results}

They apply their model to a sample of $N(0)$ atoms of ${}^7{\rm Li}$, 
with some specific parameters (see the paper for details).
They then study the evolution of the free variational parameters 
$A(\tau)$, $a(\tau)$, and $b(\tau)$ for 
$N(0)=800$ (Fig. 1), and 
$N(0)=3500$ (Fig. 2). 

In Fig. 3, they study the stability of $|\psi|^2$ as a function of space
and time under their model for one soliton.

In Fig. 4, 5 and 6, they apply the same study as the previous one for 
a sample of $N(0)=3500$ and ifferent parameters amd prove the stability 
of their model under said parameters. \\


In this study, the authors build a matter wave and show its stability under
different sets of parameters $A(\tau)$, $a(\tau)$ and $b(\tau)$, as well
as different values of the binary interaction $g_2$ 
(see Eq. (\ref{eqn:main_eq})).
Doing so, they show the possible existence of localized modes in a trapped 
BEC with three body recombination, atomic feeding and quantum fluctuation.

Moreover, they find that the initial Gaussian form from their trial function 
breaks into three trains of solitons which oscillate during propagation
and shed light on the effect of the strength of the quanbtum fluctuations
on the the period of oscillations and the spatial separation of
localized matter wave.

This stuyd will contribute in the general understandng of the dynamics 
of matter wave and will hfind its importance in BEC applications.

\newpage

\section{Conclusion and personal considerations}

In this paper, Mboumba \textit{et al.} study the stability of matter waves , in this case solitons, trapped into Bose-Einstein condensates. Their model
take into account quantum fluctuations, three body losses as well as 
atomic feeding. They design the Gross-Pitaevskii equation with the appropriate 
parameters fdor staibility anbd demonstrate their viability of the system.
From a methodological point of view, they make use of the variational method 
in order to find an appropriate test function, and derive from the Lagrangian, 
the necessary values for the stability of the system. 

This study demonstrates the role of quantum fluctuations in the dynamics and 
allows to have an insight of the spatial behaviour of the system in space, 
during propagation. 
The main goal is to gain some insight that will ater allow to control
the solitons in BEC for applications.
Nonetheless, the study is highly theoretical and the authors fail to stress,
as far as that paper is concerned, the link between some current experimental 
data and their actual treatmnent. If not possible, it would have been very good if they at least compared their data with previous results not taking into 
account quantum fluctuations or three body losses, for example.


\end{document}
